% Source: https://artofproblemsolving.com/community/c3h2507318_mathslayer4444_college_thesis
% Stable identity: aops-topic-2507318-q01
A forum post claims to prove that $x=0.\overline{9}$ satisfies $x\ge4.705$, as follows.
\begin{enumerate}
\item[1.] The series $1-1+1-1+\cdots$ has sum $\dfrac12$.
\item[2.] We have $\dfrac{x}{11}=0.\overline{09}$ and $\dfrac{10x}{11}=0.\overline{90}$, so
\[\dfrac{9x}{11}=\dfrac{10x}{11}-\dfrac{x}{11}=\dfrac{9}{10}-\dfrac{9}{100}+\dfrac{9}{1000}-\cdots.\]
\item[3.] The Cauchy--Schwarz inequality in the form $\left(\sum a_i\right)\left(\sum b_i\right)\ge\left(\sum\sqrt{a_ib_i}\right)^2$, applied to $a_i=(-1)^{i+1}$ and $b_i=(-1)^{i+1}\dfrac{9}{10^i}$ (so that $a_ib_i=\dfrac{9}{10^i}$), gives
\[\dfrac12\cdot\dfrac{9x}{11}\ge\left(\sum_{i=1}^{\infty}\sqrt{\dfrac{9}{10^i}}\right)^2=\dfrac{9}{\left(\sqrt{10}-1\right)^2},\]
hence $x\ge\dfrac{22}{\left(\sqrt{10}-1\right)^2}\approx4.705$.
\end{enumerate}
\begin{enumerate}
\item[(a)] For each of steps 1, 2 and 3, decide whether it is valid, and explain exactly what is wrong with each invalid step.
\item[(b)] Prove that $0.\overline{9}=1$.
\end{enumerate}
